\begin{abstract}
\noindent
This note documents the models used in the paper  \emph{The Environmental Macroeconomics of the Circular Economy}. It includes a planner's problem in continuous time, a corresponding decentralized market economy, and quantitative implementations in discrete time with explicit functional forms. A defining nature of the models are that the material balance principle (mass preservation and increasing entropy) is imposed for each economic activity -- production, extraction, exploration, waste collection and treatment, consumption, investment, and recycling -- with aggregate material balance principles derived from the process-level accounting. Doing so ties the use of materials (through extraction of raw materials from nature and recycled materials from waste) together with the waste generation and pollution and yields an endogenous material intensity of final goods and capital. 


The note contains four parts that deals with the different model types: The planner's problem in continuous time, a corresponding market implementation, the quantitative version of the planner's problem in discrete time with explicit functional forms, and a corresponding market version. 


In the first section, we present the general setup for the theoretical model. In this most general framework, we allow for waste generation to occur whenever final goods are used for an economic activity (representing embodied materials diffused as waste or accumulated in durable stocks that will delay diffusion) \emph{and} proportional to activities measured in tonnes as well (extraction of new resources and exploration for new resources, both in tonnes).


The appendix presents a couple of model variants including short solution characterizations that may be either relevant to explore further in the future or be useful for the numerical implementation. We also collect a list of other model extensions/adjustments that may be worth exploring later/in separate projects. 
\end{abstract}